\section{Approche proposée}
\label{sec:approche}

Cette section présente la modélisation du problème en programmation par contraintes avec MiniZinc. Deux modèles distincts ont été développés, chacun avec une approche différente pour modéliser les contraintes de ressources et les non-chevauchements.

\subsection{Variables de décision}

Les deux modèles utilisent les mêmes variables de décision :
\begin{itemize}
    \item \texttt{start\_op[o]} : Temps de début de l'opération $o$ (domaine : $[0, H]$ où $H = 1440$ minutes)
    \item \texttt{end\_op[o]} : Temps de fin de l'opération $o$ (domaine : $[0, H]$)
    \item \texttt{start\_maint[m]} : Temps de début de la maintenance $m$ (domaine : $[0, H]$)
    \item \texttt{end\_maint[m]} : Temps de fin de la maintenance $m$ (domaine : $[0, H]$)
    \item \texttt{assigned\_track[o]} : Voie assignée à l'opération $o$ (domaine : ensemble des voies disponibles)
\end{itemize}

\subsection{Modèle 1 :}

[À compléter avec la description du Modèle 1]

\subsection{Modèle 2 : Approche alternative}

\subsection{Contraintes principales}

Les deux modèles implémentent toutes les contraintes décrites dans la Section~\ref{sec:probleme}. La différence principale réside dans la modélisation des contraintes de ressources et des non-chevauchements.

Le Modèle 2 utilise une approche différente pour modéliser les contraintes de ressources. Les contraintes les plus importantes sont :

\textbf{Capacité des ressources :} Pour vérifier que les ressources ne sont jamais dépassées, le modèle vérifie la capacité à chaque instant où une opération commence. Pour chaque opération $o$ et chaque type de technicien $k$, on somme le nombre de techniciens utilisés par toutes les opérations qui sont en cours au moment où $o$ commence. Cette somme ne doit pas dépasser le nombre de techniciens disponibles $N_k$. La même logique s'applique aux équipements avec leur capacité simultanée $C_e$.

\textbf{Non-chevauchement :} Pour deux opérations qui utilisent la même voie, le modèle impose qu'elles ne se chevauchent pas. Si elles sont du même train, il faut au moins $\sigma_{op}$ minutes entre la fin de l'une et le début de l'autre. Si elles sont de trains différents, il faut au moins $\sigma_{train}$ minutes.

\textbf{Heuristiques :} Le modèle utilise les heuristiques par défaut des solveurs MiniZinc. Chuffed utilise des techniques de propagation de contraintes et de recherche locale, tandis que Gecode utilise des stratégies de branchement (first-fail, domain splitting). Aucune heuristique personnalisée n'a été implémentée.

\subsection{Fonction objectif}

L'objectif est de minimiser le retard pondéré total : $\min \sum_{m} p_m \cdot \max(0, f_m - \delta_m)$ où $p_m$ est la priorité et $\delta_m$ la date butoir de la maintenance $m$.
